\originalchapter{核素、质量与结合能}\label{ch:nuclear-mass}
\originalsection{长度与能量尺度}
一个核素由质子数 $Z$ 和中子数 $N$ 指定, 质量数为 $A=Z+N$, 记为 ${}^{A}_{Z}X$. 同位素具有相同 $Z$, 同中异位素具有相同 $N$, 同量异位素具有相同 $A$. 电子决定化学性质, 核素的寿命、核自旋和核反应却不能由化学元素名独自确定. 核尺度通常用 $1\,\mathrm{fm}=10^{-15}\,\mathrm m$, 能量用 MeV; $\hbar c\simeq197.327\,\mathrm{MeV\,fm}$ 把波长与动量连接起来.

实验上多数核的半径满足 $R\simeq r_0A^{1/3}$, 其中 $r_0$ 约为 $1.2$--$1.3\,\mathrm{fm}$. 若以均匀球估算, 核子数密度 $A/(4\pi R^3/3)=3/(4\pi r_0^3)$ 与 $A$ 几乎无关. 这与每个核子只同邻近核子强烈相互作用的饱和性质一致. 半径标度是经验近似, 不适用于把轻核、晕核和形变核全部视为同样的均匀球.

\originalsection{质量亏损和原子质量}
\begin{definition}
核的结合能为 $B=[Zm_p+Nm_n-M_{\rm nuc}]c^2$. 它是把静止核完全拆成静止自由核子所需的能量, 对稳定束缚态取正值.
\end{definition}
核质量表常给中性原子质量 $M_{\rm at}$. 若暂忽略电子总结合能, 有 $M_{\rm nuc}\simeq M_{\rm at}-Zm_e$, 因而
\begin{equation}
 B\simeq[ZM({}^{1}\mathrm H)+Nm_n-M_{\rm at}]c^2.
\end{equation}
电子数在两端相同的时候, 使用原子质量可自动抵消电子静质量. 对需要 keV 或更高精度的重元素比较, 还须处理电子结合能. $1\,\mathrm u\,c^2\simeq931.494\,\mathrm{MeV}$ 是质量差换算常数; 不能把 $A$ 本身当作精确的原子质量.

每核子结合能 $B/A$ 在中等质量核附近较大. 因此轻核聚变和重核裂变都可能释放能量, 但``结合能变大''不是反应必然发生的充分条件: 电荷、核子数、动量和角动量必须守恒, 反应还可能受到库仑势垒与相空间限制.

\originalsection{液滴模型与稳定线}
半经验质量公式把核看成带电、近乎不可压缩的液滴:
\begin{equation}\label{eq:semf}
 B=a_vA-a_sA^{2/3}-a_c\frac{Z(Z-1)}{A^{1/3}}-a_a\frac{(A-2Z)^2}{A}+\delta(A,Z).
\end{equation}
体积项表示邻近核子的吸引; 表面核子邻居较少, 所以表面项减少结合. 均匀带电球的静电能约为 $3Z^2e^2/(20\pi\epsilon_0R)$, 解释库仑项的 $Z^2/A^{1/3}$ 标度. 不对称项源于费米填充: 质子和中子分属两种费米体系, 数量越不均衡, 更多粒子被迫占据较高动量态. 配对项在偶偶核取正、奇奇核取负、奇 $A$ 常取零. 各项系数依拟合区间改变; 作业指定系数时必须使用题目版本, 特别注意不对称项是否含 $1/4$.

固定 $A$ 先忽略配对项并用 $Z^2$ 近似 $Z(Z-1)$. 极小化原子质量
$M c^2=(A-Z)m_nc^2+Z(m_p+m_e)c^2-B$ 得
\begin{equation}
 Z_* =\frac{4a_a+(m_n-m_p-m_e)c^2}{8a_a/A+2a_c/A^{1/3}}.
\end{equation}
轻核的库仑项小, $Z_*\simeq A/2$; 重核则倾向中子较多. 实际 $Z$ 为整数, 偶 $A$ 还有两条配对质量抛物线, 需要比较附近整数而非简单四舍五入连续结果.

\begin{center}
\begin{tikzpicture}[x=0.035cm,y=0.055cm]
\draw[->] (0,0)--(240,0) node[right]{$A$};\draw[->](0,0)--(0,120)node[above]{$Z$};
\draw[dashed,gray](0,0)--(220,110)node[above]{$Z=A/2$};
\draw[domain=2:220,samples=100,thick,structurecolor]plot(\x,{\x/(2+0.015*\x^(2/3))});
\node at (170,50){液滴稳定线};
\end{tikzpicture}
\end{center}
图用示意系数呈现趋势, 不是核素质量表. 核壳结构造成的魔数、局部形变和同质异能态需要额外的微观模型.

\originalsection{分离能与反应能}
单中子分离能 $S_n=[M(A-1,Z)+m_n-M(A,Z)]c^2$ 等于移走一个中子所需能量; 它比较的是邻近核, 因而能显示 $B/A$ 的平滑曲线看不出的壳层间断. 两中子分离能 $S_{2n}$ 常能减弱奇偶配对的摆动. 对反应 $a+A\to b+B$, $Q=(m_a+m_A-m_b-m_B)c^2$. $Q>0$ 意味静质量可转为动能, $Q<0$ 则需要入射动能. 若靶静止, 出射系统质心仍须运动, 所以吸能反应阈值大于 $|Q|$.

反应能与势垒高度承担不同任务. 例如 $\alpha$ 衰变可能有正 $Q$, 但 $\alpha$ 粒子仍要隧穿核外库仑势垒, 半衰期因此可能很长. 下一章用概率流区分``存在波函数''和``真的发生传播''.
